\subsection*{6.2 Runtime of Algorithm 3}
In this subsection, we bound the runtime of Algorithm 3 and set $s$ and $t$ to the optimal values. We analyze Algorithm 3 in its parts.
\paragraph{Part 1}
The for loop in part 1 clearly runs $t$ times, so we just need to bound the running time of each iteration. Generating a uniformly random binary with $L(f)$-leaves can be done in linear time (see~\cite{ref15} for a survey of various approaches). The other operations in the for loop can clearly be done in time $O(N+L(f))$. Hence, all of part 1 runs in time $O(t\cdot(N+L(f)))$ which is $O(t\cdot N)$ using the worst-case formula upper bound from Theorem 13.
Moreover, part 1 only makes oracle calls of length $N$ (to calculate $L(f)$).
\paragraph{Part 2}
Building the $\operatorname{SmallFuncs}$ set requires iterating through all formulas of size $s$ (which is bounded by $2^{(1+o(1))\cdot s\log n}$ using Proposition 12) and then computing the truth table of each of these size $s$ formulas (which can be done in time $O(Ns)$). Hence, computing $\operatorname{SmallFuncs}$ can be done in $O(N\cdot 2^{(1+o(1))\cdot s\log n})$ time. Moreover, $|\operatorname{SmallFuncs}|$ is clearly upper bounded by the upper bound on the number of formulas of size $s$: $2^{(1+o(1))\cdot s\log n}$.
Next, building the $\Candidates$ set can be done in time $O(|\operatorname{SmallFuncs}|+N)=O(N\cdot 2^{(1+o(1))\cdot s\log n})$, and we use oracle calls of length $2N$ in this step (for $\Select[f,g]$).
Hence, part 2 runs in time $O(N\cdot 2^{(1+o(1))\cdot s\log n})$.
We will make use of the following claim later, which bounds the size of the $\Candidates$ set if part 1 did not return a value.
\begin{ntclaim}[Author Claim 51]
Fix some function $f$. Then with 99\% probability (over the algorithm's choice of random formulas) either Algorithm 3 on input $f$ returns before reaching part two or
\[|\Candidates|\leq\frac{2^{N(1+o(1))}}t.\]
\end{ntclaim}
\begin{proof}
Suppose that Algorithm 3 reaches part two on input $f$ and $|\Candidates|>2^{N(1+o(1))}/t$. Then Claim 45 implies that
\[\frac{2^{N(1+o(1))}}t=O(|\Canon{0}(f)|N\log n)\]
which implies that
\[|\Canon{0}(f)|\geq\frac{2^{N(1+o(1))}}t.\]
Hence, Claim 43 implies that Algorithm 3 will return in part 1 with 99\% probability.
\end{proof}
\paragraph{Part 3}
It is easy to see that part 3 runs in time $O(|\Candidates|^2+N)$ and makes oracle calls of length $N$.
Thus, using Claim 51, we have that with 99\% probability part 3 runs in time $2^{2N(1+o(1))}/t^2$.
\paragraph{Part 4}
Computing $\operatorname{SmallCandidates}$ can be done in $O(|\Candidates|+N)$ time, and the outer for loop runs at most $O(|\Candidates|)$ many times.
It remains to bound the running time of each iteration of the outer for loop. The if condition can be checked in $O(N)$ time. Constructing $h_{?,g}$ also takes $O(N)$ time (similar to the if condition, just iterate through each input $x\in\{0,1\}^n$ and see which values of $b\in\{0,1\}$ satisfy $f(x)=g(x)Ob$). Each iteration of the inner for loop takes $O(N)$ time. Finally, the inner for loop runs $2^{|h_{?,g}^{-1}(?)|}$ many times.
Thus, the total running time for part 4 is
\[O\left(|\Candidates|\cdot\left(N+N\cdot 2^{\max_g\{|h_{?,g}^{-1}(?)|\}}\right)\right)\]
time.
Moreover, we can bound the quantity $\max_g\{|h_{?,g}^{-1}(?)|\}$ as follows.
\begin{ntclaim}[Author Claim 52]
$\max_g\{|h_{?,g}^{-1}(?)|\}\leq(1+o(1))(N-s\log n)$.
\end{ntclaim}
\begin{proof}
Since any function $h$ that agrees with $h_{?,g}$ satisfies $gOh=f$ and, we have that
\[L(f)\leq L(g)+L(h_{?,g}).\]
Since $L(g)\leq L(f)-s$ (since $g\in\operatorname{SmallCandidates}$), we have that $L(h_{?,g})\geq s$.
On the other hand, the upper bound on the formula complexity of partial functions from Theorem 16 implies that
\[L(h_{?,g})\leq(1+o(1))\left(1-\frac{|h_{?,g}^{-1}(?)|}N\right)\frac N{\log n}.\]
Hence
\[s\leq(1+o(1))\left(1-\frac{|h_{?,g}^{-1}(?)|}N\right)\frac N{\log n}\]
so
\[|h_{?,g}^{-1}(?)|\leq N-\frac{s\log n}{1+o(1)}\leq(1+o(1))(N-s\log n).\]
\end{proof}
Thus, using Claim 51, we have that with 99\% probability part 4 runs in $2^{2N(1+o(1))}/t^2$.
\[O\left(\frac{2^{N(1+o(1))}}t\cdot N\cdot 2^{(1+o(1))(N-s\log n)}\right)=\frac{2^{(1+o(1))(2N-s\log n)}}t\]
time.
\paragraph{In total}
Thus, we get that with 99\% probability Algorithm 3 runs in time
\[O(t\cdot N)+O(N\cdot 2^{(1+o(1))\cdot s\log n})+\frac{2^{2N(1+o(1))}}{t^2}+\frac{2^{(1+o(1))(2N-s\log n)}}t.\]
Letting $s=\frac 23\frac{2^n}{\log n}$ and $t=2^{(2/3)N}$, we get that the running time is bounded by
\[2^{(1+o(1))(2/3)N}.\]
Moreover, $s$ will satisfy $s\geq L(f)/2$ (as required for the correctness of the algorithm) for all $f$ with $n$-inputs when $n$ is sufficiently large by Theorem 13.
\end{proof}
